\documentclass[12pt]{article}
%%---------------------------------------------------------------------
% packages
% geometry
\usepackage{geometry}
% font
\usepackage{fontspec}
\defaultfontfeatures{Mapping=tex-text}  %%如果没有它，会有一些 tex 特殊字符无法正常使用，比如连字符。
\usepackage{xunicode,xltxtra}
\usepackage[BoldFont,SlantFont,CJKnumber,CJKchecksingle]{xeCJK}  % \CJKnumber{12345}: 一万二千三百四十五
\usepackage{CJKfntef}  %%实现对汉字加点、下划线等。
\usepackage{pifont}  % \ding{}
% math
\usepackage{amsmath,amsfonts,amssymb}
% color
\usepackage{color}
\usepackage{xcolor}
\definecolor{EYE}{RGB}{199,237,204}
\definecolor{FLY}{RGB}{128,0,128}
\definecolor{ZHY}{RGB}{139,0,255}
% graphics
\usepackage[americaninductors,europeanresistors]{circuitikz}
\usepackage{tikz}
\usetikzlibrary{positioning,arrows,shadows,shapes,calc,mindmap,trees,backgrounds}  % placements=positioning
\usepackage{graphicx}  % \includegraphics[]{}
\usepackage{subfigure}  %%图形或表格并排排列
% table
\usepackage{colortbl,dcolumn}  %% 彩色表格
\usepackage{multirow}
\usepackage{multicol}
\usepackage{booktabs}
% code
\usepackage{fancyvrb}
\usepackage{listings}
% title
\usepackage{titlesec}
% head/foot
\usepackage{fancyhdr}
% ref
\usepackage{hyperref}
% pagecolor
\usepackage[pagecolor={EYE}]{pagecolor}
% tightly-packed lists
\usepackage{mdwlist}

\usepackage{styles/iplouccfg}
\usepackage{styles/zhfontcfg}
\usepackage{styles/iplouclistings}

%%---------------------------------------------------------------------
% settings
% geometry
\geometry{left=2cm,right=1cm,top=2cm,bottom=2cm}  %设置 上、左、下、右 页边距
\linespread{1.5} %行间距
% font
\setCJKmainfont{Adobe Kaiti Std}
%\setmainfont[BoldFont=Adobe Garamond Pro Bold]{Apple Garamond}  % 英文字体
%\setmainfont[BoldFont=Adobe Garamond Pro Bold,SmallCapsFont=Apple Garamond,SmallCapsFeatures={Scale=0.7}]{Apple Garamond}  %%苹果字体没有SmallCaps
\setCJKmonofont{Adobe Fangsong Std}
% graphics
\graphicspath{{figures/}}
\tikzset{
    % Define standard arrow tip
    >=stealth',
    % Define style for boxes
    punkt/.style={
           rectangle,
           rounded corners,
           draw=black, very thick,
           text width=6.5em,
           minimum height=2em,
           text centered},
    % Define arrow style
    pil/.style={
           ->,
           thick,
           shorten <=2pt,
           shorten >=2pt,},
    % Define style for FlyZhyBall
    FlyZhyBall/.style={
      circle,
      minimum size=6mm,
      inner sep=0.5pt,
      ball color=red!50!blue,
      text=white,},
    % Define style for FlyZhyRectangle
    FlyZhyRectangle/.style={
      rectangle,
      rounded corners,
      minimum size=6mm,
      ball color=red!50!blue,
      text=white,},
    % Define style for zhyfly
    zhyfly/.style={
      rectangle,
      rounded corners,
      minimum size=6mm,
      ball color=red!25!blue,
      text=white,},
    % Define style for new rectangle
    nrectangle/.style={
      rectangle,
      draw=#1!50,
      fill=#1!20,
      minimum size=5mm,
      inner sep=0.1pt,}
}
\ctikzset{
  bipoles/length=.8cm
}
% code
\lstnewenvironment{VHDLcode}[1][]{%
  \lstset{
    basicstyle=\footnotesize\ttfamily\color{black},%
    columns=flexible,%
    framexleftmargin=.7mm,frame=shadowbox,%
    rulesepcolor=\color{blue},%
%    frame=single,%
    backgroundcolor=\color{yellow!20},%
    xleftmargin=1.2\fboxsep,%
    xrightmargin=.7\fboxsep,%
    numbers=left,numberstyle=\tiny\color{blue},%
    numberblanklines=false,numbersep=7pt,%
    language=VHDL%
    }\lstset{#1}}{}
\lstnewenvironment{VHDLmiddle}[1][]{%
  \lstset{
    basicstyle=\scriptsize\ttfamily\color{black},%
    columns=flexible,%
    framexleftmargin=.7mm,frame=shadowbox,%
    rulesepcolor=\color{blue},%
%    frame=single,%
    backgroundcolor=\color{yellow!20},%
    xleftmargin=1.2\fboxsep,%
    xrightmargin=.7\fboxsep,%
    numbers=left,numberstyle=\tiny\color{blue},%
    numberblanklines=false,numbersep=7pt,%
    language=VHDL%
    }\lstset{#1}}{}
\lstnewenvironment{VHDLsmall}[1][]{%
  \lstset{
    basicstyle=\tiny\ttfamily\color{black},%
    columns=flexible,%
    framexleftmargin=.7mm,frame=shadowbox,%
    rulesepcolor=\color{blue},%
%    frame=single,%
    backgroundcolor=\color{yellow!20},%
    xleftmargin=1.2\fboxsep,%
    xrightmargin=.7\fboxsep,%
    numbers=left,numberstyle=\tiny\color{blue},%
    numberblanklines=false,numbersep=7pt,%
    language=VHDL%
    }\lstset{#1}}{}
% pdf
\hypersetup{pdfpagemode=FullScreen,%
            pdfauthor={Haiyong Zheng},%
            pdftitle={Title},%
            CJKbookmarks=true,%
            bookmarksnumbered=true,%
            bookmarksopen=false,%
            plainpages=false,%
            colorlinks=true,%
            citecolor=green,%
            filecolor=magenta,%
            linkcolor=cyan,%red(default)
            urlcolor=cyan}
% section
%http://tex.stackexchange.com/questions/34288/how-to-place-a-shaded-box-around-a-section-label-and-name
\newcommand\titlebar{%
\tikz[baseline,trim left=3.1cm,trim right=3cm] {
    \fill [cyan!25] (2.5cm,-1ex) rectangle (\textwidth+3.1cm,2.5ex);
    \node [
        fill=cyan!60!white,
        anchor= base east,
        rounded rectangle,
        minimum height=3.5ex] at (3cm,0) {
        \textbf{\thesection.}
    };
}%
}
\titleformat{\section}{\Large\bf\color{blue}}{\titlebar}{0.1cm}{}
% head/foot
\setlength{\headheight}{15pt}
\pagestyle{fancy}
\fancyhf{}
\chead{\color{black!50!green}寻找分叉点过程详解}
\lfoot{\color{blue!50!green}邱欣欣}
\cfoot{\color{blue!50!green}\href{http://vision.ouc.edu.cn/~zhenghaiyong}{CVBIOUC}}
\rfoot{\color{blue!50!green}$\cdot$\ \thepage\ $\cdot$}
\renewcommand{\headrulewidth}{0.4pt}
\renewcommand{\footrulewidth}{0.4pt}

%%---------------------------------------------------------------------
\begin{document}
%%---------------------------------------------------------------------
%%---------------------------------------------------------------------
% \titlepage
\title{\vspace{-2em}寻找分叉点过程详解\vspace{-0.7em}}

\author{邱欣欣}
\date{\vspace{-0.7em}2015年3月19日\vspace{-0.7em}}
%%---------------------------------------------------------------------
\maketitle\thispagestyle{fancy}
%%---------------------------------------------------------------------
%\maketitle
%\tableofcontents 

\section{简介}
文档主要介绍寻找分叉点的主要过程，这部分程序集中在 \verb|Registration/points_init.m| 中，程序输入为二值图像 \verb|bw|，输出为其分叉点。

\section{过程}
\begin{enumerate}
\item 求取图像中每个像素的8邻域内像素值为1的像素个数，并记录在 \verb|countmap| 中；提取其中值大于等于3的像素，因为只有这样的像素点才可能是分叉点。
\item 按8连通寻找连通区域（ \verb|bwlabel| 函数），并标记了每个连通区域的标签值（标签的值为1、2、num（连通区域的个数））。
\item 根据每个连通区域内的像素个数（ \verb|regionprops| 函数），将连通区域划分为简单连通区域（像素个数<6）和复杂连通区域（像素个数>=6）两种情况。
\begin{enumerate}
\item 简单连通区域形状一般为以下情况(图\ref{fig:simple})，这种情况下直接求取该连通区域的中心处（ \verb|regionprops| 函数）作为分叉点（红色像素处）：
\begin{figure}[ht]
\centering
\includegraphics[width=0.9\textwidth]{images/simple}
\caption{简单情况}
\label{fig:simple}
\end{figure}

\item 当该连通区域像素数为6时，在多种情况中先判断是否为图$a$情况(图\ref{fig:complex1})：否则则判断是否为图$c$情况，其他则为图$b$等情况（$b$和$c$的判断与处理在\verb|specialpoints_init|函数中实现）。图$a$处理方法与简单连通区域方法相同，只选择中心处作为分叉点（考虑到实际中应该只有一个分叉点）。而图$b$、$c$及其他的情况下，需要进行连通区域的进一步分离，但两者处理方法不同：图$c$情况下需要求得该连通区域所有点的横纵坐标，直接确定2个分叉点的序号；图$b$情况则通过连通区域的像素数按像素序号简单分成两部分（图$c$不能这样处理的原因是对c的情况处理会有错误）。图$d$则是一些其他可能出现的情况（这种情况下图$d$与图$b$处理情况相同）。
\begin{figure}[ht]
\centering
\subfigure[]{
\includegraphics[width=0.26\textwidth]{images/complex1}}
\hspace{0.2in}
\subfigure[]{
\includegraphics[width=0.45\textwidth]{images/complex6}}
\subfigure[]{
\includegraphics[width=0.26\textwidth]{images/complex5}}
\hspace{0.8in}
\subfigure[]{
\includegraphics[width=0.35\textwidth]{images/complex7}}
\caption{第一步判断}
\label{fig:complex1}
\end{figure}

\item 当连通区域像素数大于6时：
\begin{enumerate}
\item 首先判断连通区域内是否有一个或以上像素其周围8邻域内有>=5个像素值为1，是则认为该连通区域表示的为多像素聚集的情况，如下面所示的几种情况(图\ref{fig:complex3})，此时与简单连通区域处理方法相同，只选择中心处作为分叉点：
\begin{figure}[ht!]
\centering
\includegraphics[width=0.3\textwidth]{images/complex3}
\caption{第二步判断}
\label{fig:complex3}
\end{figure}
\item 将剩下的连通区域划分为两部分：方法与处理像素数为6时图$b$、$c$的情况相同。如下面的几种情况(图\ref{fig:complex4})：
\begin{figure}[ht!]
\centering
\includegraphics[width=0.4\textwidth]{images/complex4}
\caption{第三步判断}
\label{fig:complex4}
\end{figure}
\end{enumerate}
\end{enumerate}
\end{enumerate}

\section{其他}
\begin{enumerate}
\item 当函数寻找的中心处像素值为0（即没有像素点存在）时，则给该点赋值为1。
\end{enumerate}
%%---------------------------------------------------------------------
\end{document}
